\chapter{Metoden finst}

\begin{motto}
At ein metode kan lærast og prøvast av andre,\\ er ikkje noko mindre enn at han er ein metode.
\end{motto}

Den hardaste skuldinga mot designfaget er ikkje at det manglar ein etikk eller eit fagspråk, men at det manglar ein metode — at det ikkje finst nokon delt, eksplisitt, overførbar framgangsmåte som skil designaren sitt arbeid frå inspirert improvisasjon. Skuldinga er hard fordi metoden, meir enn noko anna, er det som gjer eit fag til eit fag og ikkje ei samling talent. Ein kan bli ein god målar utan å kunne gjere greie for korleis ein måler; men ein kan ikkje byggje ein \emph{disiplin} på dugleik som ikkje let seg formulere, fordi ein disiplin må kunne overlevere arbeidet sitt til neste generasjon og leggje det fram for fagfellar. Difor er påstanden om «inga delt metode» det næraste kritikken kjem ved å treffe hjarteroten. Om han held, held mykje av resten av diagnosen med.

Dette kapitlet bestrider skuldinga direkte, og det gjer det ikkje ved å vise til ånd eller intuisjon, men til dokument. Design har eksplisitte metodar. Dei er skrivne ned, dei er publiserte, dei vert underviste, og fleire av dei har innebygde fasar for å validere sine eigne resultat. Den som hevdar at faget manglar ein delt metode, har ikkje lese metodelitteraturen — eller han har lese henne og avgjort på førehand at ingenting som står der, kan telje. Det fyrste er ein mangel ved lesinga; det andre er den sirkelen vi alt har avslørt: at kvar metode ein peikar på, vert omdefinert til ikkje å vere ekte metode. Uansett kva, held ikkje påstanden. Det vil seie noko meir presist: at design ikkje har éin universell metode som alle nyttar likt, er sant og uinteressant; at design ikkje har eksplisitte, delte, lærbare metodar i det heile, er usant og lett å motbevise. Dette kapitlet held seg til den andre, sterkare påstanden, og fellar henne.

\section{Då design fekk ein metode}

Det byrja ikkje med ein draum om kunst, men med eit krav om strenghet. Herbert Simon gav, i \textit{The Sciences of the Artificial}, den fyrste systematiske argumentasjonen for at det fanst ein eigen logikk for det menneskeskapte \autocite{simon1969sciences}. Hans poeng var at naturvitskapen studerer korleis verda \emph{er}, medan design — det vere seg ingeniørarbeid, medisin eller forretningsdrift — handlar om korleis verda \emph{kan bli}, om å gjere noverande tilstandar om til føretrekte. Og dette \emph{kan} studerast: det finst ein vitskap om søk gjennom rom av moglege løysingar, om dekomponering av problem, om tilfredsstillande framfor optimale val. Simon overdreiv kor langt den formelle modelleringa kunne strekkjast — det var hans tids optimisme — men han etablerte noko som varer: at design har ein indre struktur som let seg granske, ikkje berre utøve. At påstanden om «inga metode» skal kunne stå, krev at ein ser bort frå sjølve grunnverket faget byggjer på.

Det same tiåret samla ein generasjon designforskarar seg om nett dette prosjektet. Konferansen i London i 1962, dokumentert i samlinga \textit{Conference on Design Methods}, vert rekna som fødselen til designmetoderørsla \autocite{jones1963conference}; og John Christopher Jones gav henne si fyrste store utgreiing i \textit{Design Methods: Seeds of Human Futures} \autocite{jones1970methods}. Jones katalogiserte og systematiserte framgangsmåtar — for å bryte opp problem, for å utforske løysingsrommet, for å kople behov til form — og han gjorde det med ein eksplisitt ambisjon om at desse metodane skulle kunne delast, lærast og prøvast på nytt av andre. Han skilde, mellom anna, mellom den «svarte boksen» — designaren som produserer løysingar utan å kunne gjere greie for korleis — og den «glasboksen», der prosessen er gjennomsiktig og kvart steg kan granskast. Heile poenget med rørsla var å flytte design frå den fyrste mot den andre: frå udelt talent mot delbar framgangsmåte. Det er nett denne rørsla skuldinga om «inga metode» må late som ikkje fann stad.

Det er verdt å merke seg at Jones seinare vart kritisk til den meir mekaniske greina av rørsla. Dette vert ofte ført fram mot designmetoden, som om grunnleggjaren sjølv skulle ha trekt han tilbake. Men kritikken hans gjaldt overdriven systematisering — metodar så rigide at dei kvelte det menneskelege skjønet dei skulle tene — ikkje ideen om at design har metodar i det heile. Han ville ha metodar som opna prosessen, ikkje metodar som mekaniserte han. Det er ein usemje \emph{innanfor} eit fag med metode, om kva slags metode som er den rette, ikkje eit prov på at metode manglar. Tvert om: berre eit fag som \emph{har} metode kan ha ein opplyst strid om kva han skal vere. Striden sjølv er eit teikn på modning, ikkje på fråvær.

\section{Dei replikerbare metodologiane}

Den som meiner at designmetodar er for laust formulerte til å telje som ekte metode, bør møte to verk som ikkje let seg avfeie slik. Det fyrste er \textit{DRM, a Design Research Methodology} av Lucienne Blessing og Amaresh Chakrabarti \autocite{blessing2009drm}. DRM er ikkje ei samling tips; det er eit fullt rammeverk for å drive designforsking med etterprøvbar struktur. Det er bygd opp i fire etappar — ei klargjering av forskingsmålet, ei \emph{deskriptiv} studie som kartlegg den eksisterande situasjonen empirisk, ei \emph{preskriptiv} studie som utviklar støtte eller intervensjon, og endå ei deskriptiv studie som \emph{evaluerer} om intervensjonen faktisk verka. Det avgjerande her er det siste leddet. DRM har innebygde valideringsfasar: metodologien krev at forskaren spesifiserer på førehand kva som ville telje som suksess, og deretter måler om det blei nådd. Dette er ikkje retorikk om kreativ prosess. Det er ein forskingslogikk med same grunnform som den eksperimentelle: hypotese, intervensjon, måling, revisjon.

Det andre verket kjem frå informasjonssystemfaget, men gjeld design like fullt: tradisjonen for \emph{design science research}. \textcite{hevner2004dsr} formulerte sju eksplisitte retningslinjer for korleis ein skal byggje og evaluere designa artefakt på ein vitskapleg forsvarleg måte — at forskinga må produsere ein artefakt, at han må vere relevant for eit reelt problem, at han må evaluerast rigorøst, at han må yte eit klart kunnskapsbidrag, at sjølve forskinga må vere stringent, at designet må bli til gjennom eit reelt søk etter løysingar, og at resultatet må kunne formidlast til både fagfolk og praktikarar. Sju krav, kvart formulert så det kan prøvast: ein lesar kan halde eit stykke designforsking opp mot dei og avgjere kvar det held og kvar det sviktar. Det er det motsette av eit fag utan delte standardar for kva godt arbeid er.

Få år seinare la \textcite{peffers2007dsrm} til det som mangla i Hevners ramme: ein eksplisitt, stegvis prosessmodell. Deira \textit{Design Science Research Methodology} fører forskaren gjennom seks definerte aktivitetar — frå å identifisere problemet, gjennom å fastsetje mål, designe og utvikle, demonstrere og evaluere, til å kommunisere resultatet. Dei gjorde meir enn å liste stega: dei viste at ein kan gå inn i syklusen frå ulike inngangar — frå eit problem, frå eit mål, frå eit eksisterande artefakt — alt etter kva forskinga spring ut av, utan at framgangsmåten misser strukturen sin. Dette er ein protokoll, og ein fleksibel ein. Han kan følgjast, han kan brytast med, og to forskarar som bruker han på same problem kan samanlikne kva dei gjorde og kvifor utfalla skilde seg. Det er nett dette ein metode er til for: ikkje å garantere same svar, men å gjere vegen til svaret synleg og prøvbar. Å kalle dette «inga delt metode» er som å stå framfor ei oppskrift, ein protokoll og ei sjekkliste på éin gong og hevde at det ikkje finst matlaging.

\section{Mønster som overførbar kunnskap}

Det finst ein annan måte metoden tek form på, og han er kanskje den mest underkjende: mønsteret. Christopher Alexander gjorde to ting som høyrer saman her. I \textit{Notes on the Synthesis of Form} freista han fyrst det formelle — ein nær matematisk metode for å bryte eit designproblem ned i delkrav og syntetisere ei form som løyste dei \autocite{alexander1964notes}. Det var designmetoderørsla i si reinaste, mest systematiske form. Men Alexander gjorde noko meir interessant etterpå. Han forkasta den byråkratiske metoden — men ikkje kunnskapen han var meint å fange. I \textit{A Pattern Language} la han fram noko anna: eit \emph{språk} av mønster, kvart eit namngjeve, gjenkjenneleg, dokumentert grep som løyser eit tilbakevendande problem i ein kontekst \autocite{alexander1977pattern}. Eit mønster seier: i denne situasjonen, gjeve denne spenninga mellom desse kreftene, har denne forma vist seg å fungere — og her er kvifor. Det er overførbar kunnskap i si reinaste form: kondensert røynsle, abstrahert frå det einskilde tilfellet, formulert så ein annan kan bruke han og prøve om han held.

Det er viktig å lese Alexander rett her, for han vert ofte brukt mot designmetoden, som om han skulle ha mista trua på at design kan vere systematisk. Det stemmer ikkje. Det han forkasta, var førestillinga om at god form kan produserast av ein \emph{prosedyre} løyst frå menneskeleg dom — den mekaniske draumen. Det han heldt fast på, var at designkunnskap kan namngjevast, ordnast og delast. Mønsterspråket er ikkje eit fråfall frå metoden; det er ei mognare form av han. Og som det neste kapitlet skal vise, viste denne forma seg så robust at ho lét seg flytte til heilt andre fagfelt og framleis bere kunnskap. Dette er det motsette av eit fag utan delte reiskapar. Det er eit fag som fann ein måte å akkumulere innsikt på som er presis nok til å krysse domenegrenser.

\section{Den lærte metoden}

Det finst eit svar kritikaren kunne gje på alt dette, og det er verdt å ta fram i lyset. Han kunne seie: vel, det finst nedskrivne metodologiar — DRM, design science, mønsterspråk — men dei vert ikkje brukte i den daglege designpraksisen; dei er akademiske byggverk som lever i bøker, ikkje i atelier og på teikneromet. Det er ein alvorleg innvending, men han bommar, fordi han overser kvar designmetoden faktisk vert overført: i undervisninga. Sjå på korleis design vert lært. Det skjer i atelieret, gjennom kritikken, gjennom porteføljen — eit pedagogisk apparat som er bemerkelsesverdig likt over heile verda, frå Oslo til Eindhoven til Pasadena. Ein student legg fram arbeidet sitt; eit panel av lærarar og medstudentar prøver det, peikar på kvar det held og kvar det sviktar, og krev at studenten gjer greie for vala sine. Dette er ikkje formlaus rettleiing. Det er ein metode for å overføre dom — ein systematisk praksis for å gjere taus designdugleik eksplisitt nok til å kritiserast og forbetrast.

Poenget er at metode i design ikkje berre finst som tekst; han finst som \emph{praksis ein lærer}. Når ein designstudent gjennom åra lærer å rame eit problem, generere alternativ, prototype raskt, teste mot brukarar og iterere, lærer han ein metode — overført ikkje fyrst og fremst gjennom føresegner, men gjennom å gjere arbeidet under rettleiing av nokon som meistrar det, slik kirurgen lærer å operere og advokaten lærer å prosedere. At ein del av denne metoden er taus, gjer han ikkje uoverførbar; han overførast slik all profesjonskunnskap overførast, gjennom opplæring i ein praksis. Og det avgjerande er at praksisen er \emph{delt}: to designarar utdanna i kvar sin ende av verda kjenner att kvarandre sin framgangsmåte, kan kritisere kvarandre sitt arbeid med felles omgrep, og kan samarbeide utan å måtte byggje metoden frå grunnen. Eit fag som kan overlevere arbeidsmåten sin så påliteleg frå generasjon til generasjon, har per definisjon ein delt metode. Det er nett det overlevering vil seie.

\section{Kva «metode» eigentleg krev}

Lat oss vere ærlege om kva motparten kunne svare, for forsvaret er ikkje verdt noko om det ikkje toler det. Innvendinga lyder: design har riktig nok \emph{nedskrivne} metodar, men dei er ikkje delte slik fysikken sin metode er delt — to designarar bruker DRM eller mønsterspråk ulikt, og kjem til ulike resultat, så metoden tvingar ikkje fram semje. Det er sant. Men det er ikkje eit særtrekk ved design; det er eit særtrekk ved kvar metode som rettar seg mot det opne, det kontingente, det menneskelege. To klinikarar som følgjer same diagnostiske protokoll kan komme til ulik diagnose; to historikarar med same kjeldekritiske metode les kjeldene ulikt; to ingeniørar med same dimensjoneringsmetode vel ulike løysingar. Ein metode er ikkje ein maskin som produserer eitt svar. Ein metode er ein \emph{disiplinert framgangsmåte} — eit sett steg, omsyn og prøvingar som gjer arbeidet etterprøvbart for fagfellar, slik at dei kan sjå kva som blei gjort, vurdere om det blei gjort godt, og byggje vidare. I den forstanden, som er den einaste forstanden som gjeld for profesjonsfag, har design metode i overflod.

Skiljet som tel, er ikkje mellom metode og ikkje-metode. Det er mellom ein \emph{eksplisitt} metode, som er skriven ned så han kan kritiserast og forbetrast, og ein \emph{taus} dugleik som berre finst i hovudet til utøvaren. Poenget med dette kapitlet er at design har det fyrste, ikkje berre det andre. Simon, Jones, Blessing og Chakrabarti, Hevner, Peffers, Alexander — dei har alle gjort den tause designdugleiken eksplisitt, sett han på papir, gjort han til noko ein kan lære i ei undervisning og prøve i ei studie. Det finst framleis taus dugleik i god design, slik det finst i god kirurgi og god rettsprosedyre; men over den tause dugleiken ligg eit lag av eksplisitt, delt, lærbar metode. Skuldinga om at det laget ikkje finst, er rett og slett feil, og ho fell ved fyrste møte med litteraturen.

\vignett

Men ein metode er berre så god som kunnskapen han produserer. Og det neste spørsmålet er om sjølve forskinga gjennom design — det å lage noko som ein måte å vite på — kan gje kunnskap andre kan stole på og byggje vidare på. Det er emnet for neste kapittel.

\vidarelesing{\fullcite{blessing2009drm}; \fullcite{hevner2004dsr}; \fullcite{peffers2007dsrm}; \fullcite{simon1969sciences}; \fullcite{jones1970methods}; \fullcite{alexander1977pattern}.}
